\usepackage[margin=30truemm]{geometry}
\usepackage{fancyhdr,appendix,graphicx,layout}
%\usepackage{stix2}
\usepackage{amsmath,amsthm,amssymb}
\usepackage[shortlabels]{enumitem}
\setlist[enumerate]{font=\normalfont}
\usepackage{physics}

\usepackage[numbers]{natbib}
\usepackage{bbm}
\usepackage{bm}
\usepackage{tikz}
\usepackage{tikz-cd}
\usetikzlibrary{positioning,arrows.meta,decorations.markings,calc}
\tikzset{
    cells={font=\everymath\expandafter{\the\everymath\displaystyle}},
}
\usepackage{caption}
\usepackage[pagewise]{lineno}

\usepackage{float}

\usepackage{color}
\definecolor{darkgreen}{rgb}{0,0.7,0}
\newcommand{\comment}[1]{{\color{darkgreen}#1}}
\definecolor{darkblue}{rgb}{0,0,0.7} % darkblue color
\newcommand{\defn}[1]{\textsl{\color{darkblue} #1}}
\newcommand{\old}[1]{{\color[rgb]{1,0,0}#1}}
\newcommand{\new}[1]{{\color[rgb]{0,0,1}#1}}
\newcommand{\bhx}[1]{{\color[rgb]{0.9,0.7,0.2}#1}}

%\renewcommand{\emph}[1]{\textcolor{blue}{\textit{#1}}}

\usepackage[hyperpageref]{backref}
\usepackage[pagebackref]{hyperref}
\hypersetup{
    colorlinks=true,
    linkcolor=red,  %section,table of contents
    filecolor=blue,  %local file link
    urlcolor=orange, %url
    citecolor=violet   %citation
}
\renewcommand*{\backref}[1]{}  
   \renewcommand*{\backrefalt}[4]{
      \ifcase #1 
         Not cited.
      \or
         Cited on page #2.
      \else
         Cited on pages #2.
      \fi}

\usepackage{cleveref}
\newcommand{\hyt}[1]{\hypertarget{#1}{\textit{\textcolor{Blue}{#1}}}}
\newcommand{\hyl}[1]{\hyperlink{#1}{#1}}

%%%%%%%%%%%%%%%%%%%%%%%%%%%%%%%%%%%%%%%%%%%%%%%%%
\makeatletter
\newcommand\ReDeclareMathOperator[2]{%
    \begingroup \escapechar\m@ne\xdef\@gtempa{{\string#1}}\endgroup
    \expandafter\@ifundefined\@gtempa
        {\@latex@error{Command \string#1 undefined}\@ehc}
        \relax
    \let\@ifdefinable\@rc@ifdefinable
    \DeclareMathOperator#1{#2}}
\makeatother

%%%%%%%%%%%%%%%%%%%%%%%%%%%%%%%%%%%%%%%%%%%%%%%%%%%%

% Crefname
\crefname{thm}{Theorem}{Theorems}
\crefname{dfn}{Definition}{Definitions}
\crefname{dfnprop}{Definition-Proposition}{Definition-Proposition}
\crefname{prop}{Proposition}{Propositions}
\crefname{lem}{Lemma}{Lemmas}
\crefname{cor}{Corollary}{Corollaries}
\crefname{clm}{Claim}{Claims}
\crefname{ass}{Assumption}{Assumption}
\crefname{cond}{Condition}{Condition}
\crefname{nota}{Notation}{Notation}
\crefname{conj}{Conjecture}{Conjecture}
\crefname{fct}{Fact}{Facts}
\crefname{rmk}{Remark}{Remarks}
\crefname{eg}{Example}{Examples}
\crefname{que}{Question}{Questions}
\crefname{figure}{Figure}{Figures}
\crefname{table}{Table}{Tables}
\crefname{section}{Section}{Sections}
\crefname{subsection}{Subsection}{Subsections}
\crefname{appendix}{Appendix}{Appendices}
\crefname{equation}{}{}
\crefname{main}{Theorem}{Theorems}

% Theoremstyle
\theoremstyle{plain}
\newtheorem{thm}{Theorem}[section]
\newtheorem{main}{Theorem}
\newtheorem{prop}[thm]{Proposition}
\newtheorem{lem}[thm]{Lemma}
\newtheorem{cor}[thm]{Corollary}

\theoremstyle{definition}
\newtheorem{dfn}[thm]{Definition}
\newtheorem{dfnprop}[thm]{Definition-Proposition}
\newtheorem{dfnlem}[thm]{Definition-Lemma}
\newtheorem{clm}[thm]{Claim}
\newtheorem{ass}[thm]{Assumption}
\newtheorem{cond}[thm]{Condition}
\newtheorem{nota}[thm]{Notation}
\newtheorem{conj}[thm]{Conjecture}
\newtheorem{fct}[thm]{Fact}
\newtheorem{eg}[thm]{Example}
\newtheorem{que}[thm]{Question}

\theoremstyle{remark}
\newtheorem{rmk}[thm]{Remark}
\numberwithin{equation}{section}
\makeatletter
\let\c@equation\c@thm
\makeatother
\renewcommand{\themain}{\Alph{main}}

% Letter
\newcommand{\A}{\mathcal A}
\newcommand{\D}{\mathcal D}
\renewcommand{\H}{\mathcal H}
\newcommand{\T}{\mathcal T}
\newcommand{\U}{\mathcal U}
\newcommand{\V}{\mathcal V}
\newcommand{\W}{\mathcal W}

\newcommand{\RR}{\mathbf{R}}
\newcommand{\LL}{\mathbf{L}}

\newcommand{\NN}{\mathbb N}
\newcommand{\ZZ}{\mathbb{Z}}


%Arrow
\newcommand{\mono}{\hookrightarrow}
\newcommand{\epi}{\twoheadrightarrow}
\newcommand{\iso}{\similarrightarrow}
\newcommand{\xto}{\xrightarrow}
\newcommand{\LR}{\Leftrightarrow}
\newcommand{\To}{\Rightarrow}


% Abelian Category
\DeclareMathOperator{\Fac}{Fac}
\DeclareMathOperator{\Sub}{Sub}
\DeclareMathOperator{\Filt}{Filt}
\DeclareMathOperator{\add}{add}
\DeclareMathOperator{\Add}{Add}
\DeclareMathOperator{\rad}{rad}
\ReDeclareMathOperator{\top}{top}
\DeclareMathOperator{\simp}{Sim}
\DeclareMathOperator{\id}{id}
\DeclareMathOperator{\ind}{ind}
\DeclareMathOperator{\cok}{cok}
\DeclareMathOperator{\proj}{proj}
\DeclareMathOperator{\inj}{inj}
\DeclareMathOperator{\colim}{colim}
\DeclareMathOperator{\obj}{obj}
\DeclareMathOperator{\Mod}{\mathcal Mod}
\let\mod\relax
\DeclareMathOperator{\mod}{mod}
\DeclareMathOperator{\fl}{fl}
\DeclareMathOperator{\grmod}{grmod}
\DeclareMathOperator{\Grmod}{Grmod}
\DeclareMathOperator{\sdim}{\st{dim}}



%Triangulated category
\DeclareMathOperator{\thick}{thick}
\DeclareMathOperator{\cone}{cone}
\DeclareMathOperator{\per}{per}
\DeclareMathOperator{\pvd}{pvd}


%Functor
\DeclareMathOperator{\Hom}{Hom}
\DeclareMathOperator{\End}{End}
\DeclareMathOperator{\Aut}{Aut}
\DeclareMathOperator{\RHom}{\RR Hom}
\DeclareMathOperator{\REnd}{\RR End}
\DeclareMathOperator*{\ten}{\otimes}
\DeclareMathOperator*{\Ltensor}{\otimes^{\LL}}
\DeclareMathOperator{\Tot}{Tot}
\DeclareMathOperator{\Ext}{Ext}
\DeclareMathOperator{\Tor}{Tor}


% Non-crossing partition
\DeclareMathOperator{\NC}{NC}
\DeclareMathOperator{\wide}{wide}
\DeclareMathOperator{\Kr}{Kr}
\DeclareMathOperator{\cox}{cox}
\ReDeclareMathOperator{\l}{\ell}
\newcommand{\ora}{\overrightarrow}

% Others
\newcommand{\st}[1]{\underline{#1}}
\newcommand{\ts}[1]{\overline{#1}}
\DeclareMathOperator{\la}{\langle}
\DeclareMathOperator{\ra}{\rangle}
\newcommand{\wt}{\widetilde}

\newcommand{\val}{\operatorname{val}}
\newcommand{\Max}{\operatorname{Max}}
\newcommand{\Ker}{\operatorname{Ker}}
\newcommand{\Image}{\operatorname{Im}}
\newcommand{\vv}{\mathsf{v}}
\newcommand{\ww}{\mathsf{w}}

\newcommand{\BOne}{\left(\begin{smallmatrix} 1&0\\ 0&0\\ 0&0 \end{smallmatrix}\right)}
\newcommand{\AOne}{\left(\begin{smallmatrix} 0&0&0\\ 0&0&1 \end{smallmatrix}\right)}

\newcommand{\BTwo}{\left(\begin{smallmatrix} 0&0\\ 0&0\\ 1&0 \end{smallmatrix}\right)}
\newcommand{\ATwo}{\left(\begin{smallmatrix} 0&0&0\\ 1&0&0 \end{smallmatrix}\right)}

\newcommand{\BThree}{\left(\begin{smallmatrix} 1&0\\ 1&0\\ 1&0 \end{smallmatrix}\right)}
\newcommand{\AThree}{\left(\begin{smallmatrix} 0&0&0\\ -\lambda&1&\lambda-1 \end{smallmatrix}\right)}

\newcommand{\BOneN}{%
\left(\begin{smallmatrix}
I_n&0\\
0&0\\
0&0
\end{smallmatrix}\right)}

\newcommand{\AOneN}{%
\left(\begin{smallmatrix}
0&0&0\\
0&0&I_n
\end{smallmatrix}\right)}

\newcommand{\BTwoN}{%
\left(\begin{smallmatrix}
0&0\\
0&0\\
I_n&0
\end{smallmatrix}\right)}

\newcommand{\ATwoN}{%
\left(\begin{smallmatrix}
0&0&0\\
I_n&0&0
\end{smallmatrix}\right)}

\newcommand{\BThreeN}{%
\left(\begin{smallmatrix}
I_n&0\\
I_n&0\\
I_n&0
\end{smallmatrix}\right)}

\newcommand{\AThreeN}{%
\left(\begin{smallmatrix}
0&0&0\\
-J_n&I_n&J_n-I_n
\end{smallmatrix}\right)}

\newcommand{\morsethickarrow}{%
  \mathrel{\tikz[baseline=-0.55ex]
    \draw[->] (0,0) -- (1.6em,0);}}
\newcommand{\morsedottedarrow}{%
  \mathrel{\tikz[baseline=-0.55ex]
    \draw[densely dotted,line cap=round,->]
    (0,0) -- (1.6em,0);}}

\DeclareMathOperator{\domdim}{domdim}
\DeclareMathOperator{\fdim}{fdim}
\DeclareMathOperator{\idim}{idim}
\DeclareMathOperator{\pdim}{pdim}
\DeclareMathOperator{\gldim}{gldim}
\DeclareMathOperator{\supp}{supp}
\DeclareMathOperator{\Crit}{Crit}
\DeclareMathOperator{\Ind}{Ind}
\DeclareMathOperator{\sd}{sd}
\DeclareMathOperator{\weight}{wt}
\DeclareMathOperator{\diam}{diam}
\newcommand{\SDiff}{\mathbin{\triangle}}
\DeclareMathOperator{\MH}{MH}
\DeclareMathOperator{\lat}{lat}
